% Input-ready conventional-proof component.
% Requires amsmath, amssymb, amsthm and theorem/lemma/proposition/corollary environments.
% Labels and equation labels in this file have the prefix mf:.
\providecommand{\per}{\operatorname{per}}

\section{Finite-rank mean-field perturbations}\label{mf:setup}

The next results give an asymptotic criterion for a negative endpoint
derivative and a strict obstruction in the subcritical regime. They are
conventional mathematical proofs. Their measure and limiting arguments
are not asserted to form a completed Lean theorem.

Fix a finite integer \(r\ge1\), a real number \(t>0\), and ordered vectors
\(v_{i,N}\in\mathbb R^r\), \(1\le i\le N\), with
\(\|v_{i,N}\|\le L\) for a constant \(L<\infty\) independent of \(N\).
Write \(V_N\) for the matrix with these rows and set
\begin{equation}\label{mf:matrices}
 B_N=V_NV_N^T,\qquad
 A_N=I_N+\frac tN B_N,\qquad
 \Sigma_N=\frac1N\sum_{i=1}^N v_{i,N}v_{i,N}^T.
\end{equation}
Every \(A_N\) is real positive definite. Suppose initially that the ordered
empirical measures satisfy
\begin{equation}\label{mf:empirical}
 \mu_N=\frac1N\sum_{i=1}^N\delta_{(i/N,v_{i,N})}
 \ \Longrightarrow\ \mu
\end{equation}
weakly on \([0,1]\times\overline B(0,L)\). The first marginal of \(\mu\)
is Lebesgue measure, and
\[
 \Sigma=\int vv^T\,d\mu(x,v)=\lim_N\Sigma_N.
\]
The order coordinate \(x\) is retained throughout.

For an ordered matrix \(A\), let \(D_1(A)=P'_1(A)\), with the
inversion-weighted \(q\)-permanent as defined in the main text.
For Gram rows \(w_i\), put
\[
 f(z)=\prod_i(w_i\cdot z),\qquad
 G(z)=\sum_{i<j}(w_i\wedge w_j)
                    \prod_{a\ne i,j}(w_a\cdot z),\qquad
 T(A)=\|G\|_{\mathrm F}^2.
\]
Here the Fischer norm is
\(\|\sum_\alpha c_\alpha z^\alpha\|_{\mathrm F}^2
=\sum_\alpha\alpha!|c_\alpha|^2\), summed over exterior components.
The identities
\begin{equation}\label{mf:endpoint}
 \per A=\|f\|_{\mathrm F}^2,\qquad
 2D_1(A)=\binom N2\per A-T(A)
\end{equation}
follow by expanding the Fischer inner products: the coefficient of a
permutation monomial in \(T(A)\) is
\(\binom N2-2\operatorname{inv}(\sigma)\).
The ordered wedge identity is also used in Kitamura's public proof
package \cite{Kitamura2026}.
Consequently, for positive permanent,
\begin{equation}\label{mf:ratio}
 \frac{D_1(A)}{N(N-1)\per A}
 =\frac14\left(1-\frac{T(A)}{\binom N2\per A}\right).
\end{equation}

\subsection{The general endpoint criterion}

For \(h\in\mathbb C^r\), define
\begin{align}
 \Phi(h)&=\int\log(1+t|v\cdot h|^2)\,d\mu(x,v)-\|h\|^2,
       &K&=\operatorname*{arg\,max}_{h\in\mathbb C^r}\Phi(h),\label{mf:potential}\\
 \ell_h(v)&=v\cdot h,
       &p_h(v)&=\frac{t|\ell_h(v)|^2}{1+t|\ell_h(v)|^2},\label{mf:p}\\
 q_h(v)&=\frac{t\overline{\ell_h(v)}\,v}{1+t|\ell_h(v)|^2},
       &m_h&=\int p_h(v)\,d\mu(x,v),\label{mf:q}\\
 Q_h(s)&=\int_{\{x<s\}}q_h(v)\,d\mu(x,v),
       &Q_h&=\int q_h(v)\,d\mu(x,v).\label{mf:prefix}
\end{align}
The vector-valued function \(q_h\) is continuous, including at
\(\ell_h(v)=0\). Put
\begin{equation}\label{mf:ordered-wedge}
 W_h=\iint_{\{x<y\}}q_h(v)\wedge q_h(w)\,
                       d\mu(x,v)\,d\mu(y,w)
\end{equation}
and
\begin{equation}\label{mf:score}
 \mathcal R(h)=2\|W_h\|^2+
 2\int(1-p_h(v))\|Q_h-2Q_h(x)\|^2\,d\mu(x,v)
 +(1-m_h)^2.
\end{equation}

\begin{theorem}[Full mean-field endpoint criterion]\label{mf:criterion}
Under \eqref{mf:empirical}, \(K\) is nonempty and compact, and
\begin{align}
 \min_{h\in K}\mathcal R(h)
 &\le\liminf_{N\to\infty}\frac{T(A_N)}{\binom N2\per A_N}\notag\\
 &\le\limsup_{N\to\infty}\frac{T(A_N)}{\binom N2\per A_N}
 \le\max_{h\in K}\mathcal R(h).\label{mf:criterion-bounds}
\end{align}
If \(\min_K\mathcal R>1\), then \(D_1(A_N)<0\) for all sufficiently
large \(N\). If \(\mathcal R=\mathcal R_*\) on \(K\), then
\begin{equation}\label{mf:criterion-limit}
 \frac{D_1(A_N)}{N(N-1)\per A_N}\longrightarrow
                  \frac{1-\mathcal R_*}{4}.
\end{equation}
At every \(h\in K\),
\begin{equation}\label{mf:stationarity}
 Q_h=\overline h,\qquad m_h=\|h\|^2<1.
\end{equation}
\end{theorem}

Both \(\Phi\) and \(\mathcal R\) are invariant under a global phase of
\(h\), and under conjugation because the rows are real.
Thus two conjugate phase circles of maxima give a single limiting score.
A simpler sufficient criterion is
\(\min_{h\in K}2\|W_h\|^2>1\). The other two terms in
\eqref{mf:score} are nonnegative.

\begin{corollary}[Product ordered limits]\label{mf:product-limit}
Suppose that the ordered empirical limit is
\(\mu(dx,dv)=dx\,\nu(dv)\), so that its row law does not depend on
the position coordinate. Then, for every \(h\in K\),
\begin{equation}\label{mf:product-score}
 \mathcal R(h)=1-\frac{m_h(4-m_h)}3.
\end{equation}
If \(t\lambda_{\max}(\Sigma)>1\), then
\begin{equation}\label{mf:product-positive}
 \liminf_{N\to\infty}
 \frac{D_1(A_N)}{N(N-1)\per A_N}
 \ge\min_{h\in K}\frac{m_h(4-m_h)}{12}>0.
\end{equation}
In particular the endpoint derivative is eventually positive.
\end{corollary}
\begin{proof}
The product form of \(\mu\) gives \(Q_h(s)=sQ_h\) and
\(W_h=\tfrac12 Q_h\wedge Q_h=0\). At a maximizer,
\eqref{mf:stationarity} gives \(\|Q_h\|^2=m_h\). Since \(p_h\)
also depends only on the row, the middle term of \eqref{mf:score} is
\[
 2(1-m_h)\|Q_h\|^2\int_0^1(1-2s)^2\,ds
 =\frac23m_h(1-m_h).
\]
Adding \((1-m_h)^2\) proves \eqref{mf:product-score}.

Let \(u\) be a real unit top eigenvector of \(\Sigma\). The elementary
bound \(\log(1+s)\ge s-s^2/2\) for \(s\ge0\) and the row bound
\(\|v\|\le L\) yield
\[
 \Phi(\epsilon u)
 \ge\epsilon^2\bigl(t\lambda_{\max}(\Sigma)-1\bigr)
       -\frac{t^2L^4}{2}\epsilon^4>0
\]
for sufficiently small \(\epsilon>0\). Since \(\Phi(0)=0\), zero
does not belong to \(K\). Compactness of \(K\) and
\(m_h=\|h\|^2<1\) therefore imply
\(\min_K m_h(4-m_h)>0\).
Combining \eqref{mf:product-score}, \eqref{mf:criterion-bounds}, and
\eqref{mf:ratio} proves \eqref{mf:product-positive}.
\end{proof}

The hypothesis is a product form of the ordered empirical limit itself;
no probabilistic assertion about finite random row orders is used.

\subsection{A strict obstruction below the covariance threshold}

\begin{theorem}[Subcritical endpoint increase]\label{mf:subcritical}
Suppose \(\Sigma\ne0\) and \(t\lambda_{\max}(\Sigma)<1\).
Define
\[
 R_0=(I_r-t\Sigma)^{-1},\qquad
 M_s=\int_{\{x<s\}}vv^T\,d\mu(x,v).
\]
Then
\begin{equation}\label{mf:permanent-limit}
 \per A_N\longrightarrow\det(I_r-t\Sigma)^{-1}
\end{equation}
and
\begin{equation}\label{mf:subcritical-limit}
 \frac{D_1(A_N)}{N\per A_N}\longrightarrow
 2t^2\int_0^1\operatorname{tr}
       \bigl(M_sR_0(\Sigma-M_s)\bigr)\,ds>0.
\end{equation}
In particular \(D_1(A_N)>0\) for all sufficiently large \(N\).
\end{theorem}

\begin{corollary}[Arbitrary changing orders]\label{mf:arbitrary-orders}
The ordered weak-convergence assumption can be dropped from the strict
sign conclusion. If only \(\Sigma_N\to\Sigma\ne0\) is assumed, with the
same fixed rank and row bound, and \(t\lambda_{\max}(\Sigma)<1\), then
\begin{equation}\label{mf:uniform-positive}
 \liminf_N\frac{D_1(A_N)}{N\per A_N}
 \ge\frac{t^2(\operatorname{tr}\Sigma^2)^2}{8L^4}>0.
\end{equation}
More generally, eventual bounds
\(t\lambda_{\max}(\Sigma_N)\le1-\varepsilon\) and
\(\operatorname{tr}(\Sigma_N^2)\ge s>0\), for fixed
\(\varepsilon,s>0\), imply the lower bound \(t^2s^2/(8L^4)\).
\end{corollary}

For \(N>r\), the least eigenvalue of \(A_N\) is one, while its largest
eigenvalue is \(1+t\lambda_{\max}(\Sigma_N)\).
Thus Theorem~\ref{mf:subcritical} gives a strict obstruction for this
entire bounded-row, fixed-rank mean-field construction below limiting
condition number two. It is not a theorem for every positive-definite
matrix of condition number below two, and does not include the critical
boundary \(t\lambda_{\max}(\Sigma)=1\).

\section{Proof of the full endpoint criterion}\label{mf:criterion-proof}

\subsection{Gaussian representation and concentration}

Gaussian representations of permanents also appear in the analysis of
spatial random permutations \cite{FyodorovMuirhead2021}.
We give the specific identities and limiting arguments needed here.
Let \(Z\) be a standard complex Gaussian vector in \(\mathbb C^r\), with
density \(\pi^{-r}e^{-\|z\|^2}\). Expanding by principal subsets and using
the Fischer identity gives
\begin{equation}\label{mf:gaussian}
 \per A_N=
 \mathbb E\prod_{i=1}^N\left(1+\frac tN|v_i\cdot Z|^2\right).
\end{equation}
After \(z=\sqrt N\,h\), the corresponding probability measure is
\begin{equation}\label{mf:tilted}
 \Pi_N(dh)=
 \frac{e^{N\Phi_N(h)}\,dh}{\int e^{N\Phi_N(u)}\,du},\qquad
 \Phi_N(h)=\frac1N\sum_i\log(1+t|v_i\cdot h|^2)-\|h\|^2.
\end{equation}
The functions \(\Phi_N\) converge to \(\Phi\) locally uniformly.
Indeed, weak convergence and uniform continuity on compact products give
uniform convergence on a finite \(h\)-net, and then on each compact set.
Uniformly in \(N\),
\begin{equation}\label{mf:coercive}
 \Phi_N(h)\le\log(1+tL^2\|h\|^2)-\|h\|^2.
\end{equation}
The right side is coercive and is at most \(-\|h\|^2/2\) outside a
fixed ball. This proves the compactness and nonemptiness of \(K\).

The measures \(\Pi_N\) concentrate in every neighborhood of \(K\).
To see this, take a large ball containing \(K\) and the region where
the upper bound in \eqref{mf:coercive} may be nonnegative.
On the compact complement of a neighborhood of \(K\), \(\Phi\) is
separated from its maximum by a strictly positive amount.
Local uniform convergence preserves a smaller gap for \(\Phi_N\).
A small ball around one maximizer gives a denominator bounded below by
its fixed volume times \(e^{N(\max\Phi-\epsilon)}\).
The outer tail is controlled uniformly by
\(e^{-N\|h\|^2/2}\). Choosing \(\epsilon\) smaller than the compact gap
proves the assertion. No nondegeneracy of the Hessian is needed.

Differentiating \(\Phi\) in the real and imaginary coordinates is
legitimate by the uniform row bound. Its stationary equation is
\(Q_h=\overline h\). Taking its bilinear scalar product with \(h\)
gives \(m_h=\|h\|^2\). Since \(p_h(v)<1\) everywhere, \(m_h<1\).
This proves \eqref{mf:stationarity}.

\subsection{An exact three-sector identity}

Suppress the index \(N\). For a retained set \(I\subset[N]\), put
\(J=I^c\), \(k=|I|\), and
\begin{align*}
 f_I(z)&=\prod_{i\in I}(v_i\cdot z),\\
 G_I(z)&=\sum_{\substack{i<j\\i,j\in I}}
        (v_i\wedge v_j)\prod_{b\in I\setminus\{i,j\}}(v_b\cdot z),\\
 H_{a,I}(z)&=\sum_{i\in I}\operatorname{sign}(a-i)\,v_i
                       \prod_{b\in I\setminus\{i\}}(v_b\cdot z),
                       \qquad a\in J.
\end{align*}
Empty products equal one and empty sums equal zero.

\begin{lemma}[All exterior sectors]\label{mf:sector-identity}
One has the exact identity
\begin{equation}\label{mf:three-sectors}
 T(A_N)=\sum_{I\subset[N]}\left(\frac tN\right)^{|I|}
 \left(
  \|G_I\|_{\mathrm F}^2+
  \sum_{a\in I^c}\|H_{a,I}\|_{\mathrm F}^2+
  \binom{|I^c|}{2}\|f_I\|_{\mathrm F}^2
 \right).
\end{equation}
\end{lemma}
\begin{proof}
Use Gram rows
\(w_i=\sqrt{t/N}\,v_i\oplus e_i\).
Split their wedge polynomial into the exterior sectors
\(\Lambda^2\mathbb C^r\), \(\mathbb C^r\wedge\mathbb C^N\), and
\(\Lambda^2\mathbb C^N\).
In the first sector, the coefficient of the added-variable monomial
\(\prod_{a\in J}x_a\) is \((t/N)^{k/2}G_I\).
In the second, with exterior index \(a\in J\), the coefficient of
\(\prod_{b\in J\setminus\{a\}}x_b\) is
\((t/N)^{k/2}H_{a,I}\), up to one irrelevant common sign.
In the third sector, each exterior pair in \(J\) gives
\((t/N)^{k/2}f_I\) multiplying the corresponding remaining monomial.
Distinct added-variable monomials are Fischer-orthogonal and each has
factorial weight one. The three exterior sectors are orthogonal.
Taking squared norms proves \eqref{mf:three-sectors}.
\end{proof}

For clarity, this also identifies precisely what diagonal dilution adds
to the derivative. Define
\[
 F(B)_{ij}=b_{ij}\per B(i\mid j),\qquad
 \operatorname{cut}_F(B[I];a)=
       \sum_{\substack{i<a<j\\i,j\in I}}F(B[I])_{ij},
\]
where retained indices keep their original labels in the cut.
The matrix \(F(B[I])\) has row sums \(\per B[I]\). Hence
\[
 \|H_{a,I}\|_{\mathrm F}^2
       =|I|\per B[I]-4\operatorname{cut}_F(B[I];a),
\]
and \eqref{mf:endpoint} gives
\begin{equation}\label{mf:derivative-dilution}
 D_1(A_N)=\sum_I\left(\frac tN\right)^{|I|}
 \left(D_1(B_N[I])
          +2\sum_{a\in I^c}\operatorname{cut}_F(B_N[I];a)\right).
\end{equation}
The cross-block cuts are nonnegative, as follows by differentiating
Lieb's cross-block scaling expansion, whose coefficients in the square
of the scaling variable are nonnegative \cite{Lieb1966}.

\subsection{Conditional Bernoulli moments}

For fixed \(z\), let \(I\) be a random retained set with independent
membership probabilities
\[
 p_i(z)=
 \frac{(t/N)|v_i\cdot z|^2}{1+(t/N)|v_i\cdot z|^2}.
\]
Summing \((t/N)^{|I|}|f_I(z)|^2\) produces exactly the product in
\eqref{mf:gaussian}. At \(z=\sqrt N\,h\), these probabilities are
\(p_h(v_i)\).

Away from the zero sets of the nonzero linear forms, let
\[
 U_i=\xi_i\,\frac{v_i}{\ell_h(v_i)},\qquad
 \xi_i\ \hbox{independent Bernoulli}(p_h(v_i)).
\]
Then
\begin{equation}\label{mf:bounded-moments}
 \mathbb E U_i=q_h(v_i),\qquad
 \mathbb E\|U_i\|^2
 =\frac{t\|v_i\|^2}{1+t|\ell_h(v_i)|^2}\le tL^2,\qquad
 \|\mathbb E U_i\|\le\frac{\sqrt t\,L}{2}.
\end{equation}
These quotient random variables are used only almost everywhere.
For each nonzero row its zero set is a complex hyperplane of Lebesgue
measure zero. Zero rows contribute zero and can be omitted from this
calculation. At a vanishing nonzero linear form, declaring \(U_i=0\)
would not reproduce the continuous limiting second moment in
\eqref{mf:bounded-moments}; that identification is not made.
Instead the moment expressions have continuous extensions in the original
polynomial integrals, and values on these null sets do not affect an
integral. The final functions \(p_h,q_h,\mathcal R\) have no quotient
singularities.

Set
\begin{align*}
 W_N(h)&=\frac1{N^2}\sum_{i<j}q_h(v_i)\wedge q_h(v_j),
 &Q_N(h)&=\frac1N\sum_iq_h(v_i),\\
 Q_{a,N}(h)&=\frac1N\sum_{i<a}q_h(v_i),
 &m_N(h)&=\frac1N\sum_ip_h(v_i).
\end{align*}
For the first exterior sector the random quotient is
\[
 \frac1N\sum_{i<j}U_i\wedge U_j,
\]
whose mean is \(NW_N\). In its variance, disjoint pairs have zero
covariance. There are \(O(N^3)\) pairs sharing an index and \(O(N^2)\)
identical pairs. Every surviving covariance is bounded by a constant
depending only on \(t,L\), using independence and
\eqref{mf:bounded-moments}. The factor \(N^{-2}\) thus makes the variance
\(O(N)\). The contribution of this sector divided by \(N^2\per A_N\)
equals
\begin{equation}\label{mf:first-sector-limit}
 \mathbb E_{\Pi_N}\|W_N(h)\|^2+O(N^{-1}).
\end{equation}

For a fixed missing index \(a\), the second-sector quotient is
\[
 \frac1{\sqrt N}\sum_{i\ne a}\operatorname{sign}(a-i)U_i.
\]
Its mean is
\(\sqrt N(2Q_{a,N}-Q_N+q_h(v_a)/N)\), and its variance is at most
\(tL^2\). Membership of \(a\) is independent of this sum. Summing over
\(a\) and using the uniform bounds gives the normalized second-sector
contribution
\begin{equation}\label{mf:second-sector-limit}
 \mathbb E_{\Pi_N}\frac1N\sum_a
 (1-p_h(v_a))\|Q_N(h)-2Q_{a,N}(h)\|^2+O(N^{-1}).
\end{equation}
The third sector counts missing pairs. Its normalized contribution is
\begin{equation}\label{mf:third-sector-limit}
 \frac12\mathbb E_{\Pi_N}(1-m_N(h))^2+O(N^{-1}).
\end{equation}
All three error estimates are uniform in \(h\).

\subsection{Convergence of the ordered functionals}

The discrete functions in
\eqref{mf:first-sector-limit}--\eqref{mf:third-sector-limit} converge to
their \(\mu\)-integral versions locally uniformly in \(h\).
For the order indicators, partition the \(x\)-interval into short
subintervals. The first marginal is Lebesgue measure, so boundaries
\(x=s\) and \(x=y\) have zero \(\mu\) and \(\mu\otimes\mu\) mass.
The uniform bound on \(q_h\) controls the contribution of the short
subintervals. This proves prefix convergence uniformly in \(s\).
A finite \(h\)-net and uniform continuity give joint local uniformity
in \(s,h\); the same argument handles the ordered double integral.

Thus \(\mathcal R\) is continuous and bounded, as are the discrete
integrands uniformly. Tightness of \(\Pi_N\), local uniform convergence,
and the three sector estimates give
\begin{equation}\label{mf:score-integral}
 \frac{T(A_N)}{\binom N2\per A_N}
       =\mathbb E_{\Pi_N}\mathcal R(h)+o(1).
\end{equation}
Concentration near \(K\) proves
\eqref{mf:criterion-bounds}. Equation~\eqref{mf:ratio} proves the
remaining assertions of Theorem~\ref{mf:criterion}.

\section{Proof of the subcritical obstruction}\label{mf:subcritical-proof}

\subsection{A dominated Gaussian limit}

Assume \(t\lambda_{\max}(\Sigma)<1\). Choose \(s>1\) and \(b<1\)
such that \(st\lambda_{\max}(\Sigma_N)\le b\) for all sufficiently
large \(N\). Since \(\log(1+x)\le x\),
\[
 \prod_i\left(1+\frac{st}{N}|v_i\cdot z|^2\right)
 \le e^{stz^*\Sigma_Nz}\le e^{b\|z\|^2}.
\]
This is integrable against the standard complex Gaussian, even after
multiplication by any fixed power of \(\|z\|\).
The bounded row assumption and the quadratic Taylor remainder show
pointwise convergence of the product with parameter \(t\) to
\(e^{tz^*\Sigma z}\). Dominated convergence gives
\[
 \per A_N\longrightarrow
 \mathbb E e^{tZ^*\Sigma Z}=\det(I-t\Sigma)^{-1},
\]
which proves \eqref{mf:permanent-limit}.

The positive principal-permanent weights
\[
 \mathbb P_N(I)=
 \frac{(t/N)^{|I|}\per B_N[I]}{\per A_N}
\]
define a random retained size \(K_N=|I|\). Its exponential moment is
\begin{equation}\label{mf:retained-mgf}
 \mathbb E s^{K_N}
 =\frac{\per(I+(st/N)B_N)}{\per A_N}\le C_s.
\end{equation}
Therefore every fixed moment of \(K_N\) is uniformly bounded.

\subsection{Polynomial bounds for compressed endpoint terms}

For any \(k\times k\) PSD Gram matrix \(B\), put \(C_k=\binom k2\).
Cauchy--Schwarz applied to the \(C_k\) wedge summands gives
\begin{align}
 T(B)
 &\le C_k\sum_{i<j}
       \det B[\{i,j\}]\per B[\{i,j\}^c]\notag\\
 &\le C_k^2\per B.\label{mf:wedge-bound}
\end{align}
The last inequality uses the elementary two-by-two inequality
\(\det\le\per\), followed by Lieb's principal-block inequality
\cite{Lieb1966}. The latter is also a consequence of the Fischer product
inequality for products of linear forms \cite{Reznick1993}.
It follows that
\begin{equation}\label{mf:derivative-bound}
 |D_1(B)|\le\frac{C_k+C_k^2}{2}\per B.
\end{equation}
The matrix \(F(B)\) is PSD: it is the Gram matrix of the vector-valued
polynomials \(v_i\prod_{j\ne i}\ell_j\).
Lieb's inequality gives
\(F(B)_{ii}=b_{ii}\per B(i\mid i)\le\per B\).
Thus \(|F(B)_{ij}|\le\per B\), and
\begin{equation}\label{mf:footrule-bound}
 \left|\sum_{i,j}|i-j|F(B)_{ij}\right|\le k^3\per B.
\end{equation}
The polynomial dependence on \(k\) is needed here; a factorial bound
would not suffice for \eqref{mf:retained-mgf}.

\subsection{Replacing inversions by the footrule}

Define
\begin{equation}\label{mf:footrule}
 \mathcal L(A)=\sum_{i,j}|i-j|a_{ij}\per A(i\mid j)
 =\sum_\sigma\left(\sum_i|i-\sigma(i)|\right)
                         \prod_i a_{i,\sigma(i)}.
\end{equation}
If a permutation fixes an index \(a\), deleting that index reduces its
inversion count by twice the number of edges crossing that cut in one
direction. It reduces its footrule by the same amount, because each
crossing edge has its length reduced by one and the crossing counts in
the two directions agree. Therefore
\[
 \operatorname{inv}(\sigma)-\sum_i|i-\sigma(i)|
\]
is unchanged by deleting fixed indices and compressing the remaining
order.
Expanding \(A_N=I+(t/N)B_N\) gives the exact identity
\begin{equation}\label{mf:compression}
 D_1(A_N)-\mathcal L(A_N)
 =\sum_I\left(\frac tN\right)^{|I|}
       \bigl(D_1(B_N[I])-\mathcal L(B_N[I])\bigr),
\end{equation}
where the submatrix quantities use their compressed order.
Equations~\eqref{mf:retained-mgf},
\eqref{mf:derivative-bound}, and \eqref{mf:footrule-bound} imply
\begin{equation}\label{mf:bounded-defect}
 |D_1(A_N)-\mathcal L(A_N)|\le C\per A_N
\end{equation}
with \(C\) independent of \(N\).

\subsection{The limiting cross trace}

For \(i\ne j\), Gaussian expansion of the cross permanent gives
\begin{equation}\label{mf:cross-minor}
 \per A_N(i\mid j)=\frac tN\,
 \mathbb E\left[
 (v_j\cdot Z)\overline{(v_i\cdot Z)}
 \prod_{a\ne i,j}\left(1+\frac tN|v_a\cdot Z|^2\right)\right].
\end{equation}
The domination used above, with the required Gaussian moment weight,
shows uniformly in \(i,j\) that the expectation in
\eqref{mf:cross-minor}, divided by \(\per A_N\), is
\(v_j^TR_0v_i+o(1)\).
Indeed, omitting two factors costs at most
\(O(N^{-1})\|Z\|^4\) times the dominating product; the full tilted
second-moment matrix converges in the fixed dimension.
Consequently
\[
 \frac{F(A_N)_{ij}}{\per A_N}
 =\frac{t^2}{N^2}(v_i^Tv_j)(v_j^TR_0v_i)+o(N^{-2})
\]
uniformly for \(i\ne j\). Weak convergence and \eqref{mf:footrule} yield
\begin{align}
 \frac{\mathcal L(A_N)}{N\per A_N}
 &\longrightarrow
 t^2\iint |x-y|(v^Tw)(w^TR_0v)\,d\mu(x,v)\,d\mu(y,w)\notag\\
 &=2t^2\int_0^1\operatorname{tr}
             \bigl(M_sR_0(\Sigma-M_s)\bigr)\,ds.\label{mf:footrule-limit}
\end{align}
The last equality integrates the two cut orientations of \(|x-y|\).
The two traces agree and are real, since \(R_0\) is a function of
\(\Sigma\).

\begin{lemma}[A resolvent cross-trace inequality]\label{mf:trace}
If \(A,B\succeq0\), \(S=A+B\), and \(t\|S\|<1\), then
\begin{equation}\label{mf:trace-inequality}
 \operatorname{tr}\bigl(A(I-tS)^{-1}B\bigr)
       \ge\operatorname{tr}(AB)\ge0.
\end{equation}
\end{lemma}
\begin{proof}
Restrict to the range of \(S\), since both \(A\) and \(B\) annihilate
its kernel. Write
\(A=S^{1/2}CS^{1/2}\), \(B=S^{1/2}(I-C)S^{1/2}\), with
\(0\preceq C\preceq I\), and diagonalize \(S\) with eigenvalues
\(\lambda_i>0\).
If \(f\ge0\) and \(x\mapsto xf(x)\) is nondecreasing on this spectrum,
direct expansion gives
\begin{align}
 \operatorname{tr}(Af(S)B)
 &=\sum_i\lambda_i^2f(\lambda_i)
                   \bigl(C_{ii}-(C^2)_{ii}\bigr)\notag\\
 &\quad+\frac12\sum_{i,j}(\lambda_i-\lambda_j)
       \bigl(\lambda_if(\lambda_i)-\lambda_jf(\lambda_j)\bigr)
                        |C_{ij}|^2\ge0.\label{mf:trace-expansion}
\end{align}
Every term is nonnegative, because \(C^2\preceq C\).
Apply this to \(f(x)=(1-tx)^{-1}-1\);
then \(xf(x)=tx^2/(1-tx)\) is nonnegative and increasing.
This proves the first inequality in \eqref{mf:trace-inequality};
the second is the standard trace inequality for two PSD matrices.
\end{proof}

The lower comparison is strictly positive after integration:
\begin{equation}\label{mf:strict-trace}
 2\int_0^1\operatorname{tr}\bigl(M_s(\Sigma-M_s)\bigr)\,ds
 =\iint|x-y||v^Tw|^2\,d\mu\,d\mu>0.
\end{equation}
The unweighted integral of \(|v^Tw|^2\) is
\(\operatorname{tr}(\Sigma^2)>0\), while \(x=y\) has zero
\(\mu\otimes\mu\) measure.
More quantitatively, \(|v^Tw|^2\le L^4\), and the mass of
\(|x-y|<\epsilon\) is at most \(2\epsilon\).
Taking \(\epsilon=\operatorname{tr}(\Sigma^2)/(4L^4)\) gives
\begin{equation}\label{mf:strict-lower-bound}
 \iint|x-y||v^Tw|^2\,d\mu\,d\mu
 \ge\frac{(\operatorname{tr}\Sigma^2)^2}{8L^4}.
\end{equation}
Equations~\eqref{mf:bounded-defect} and \eqref{mf:footrule-limit},
followed by Lemma~\ref{mf:trace}, prove
Theorem~\ref{mf:subcritical}.

\begin{proof}[Proof of Corollary~\ref{mf:arbitrary-orders}]
The ordered empirical measures lie in a compact space of probability
measures. Extract a weakly convergent subsequence from any proposed
violating subsequence. Its first marginal is Lebesgue measure and its
covariance is the corresponding limit of \(\Sigma_N\).
Apply Theorem~\ref{mf:subcritical} and
\eqref{mf:strict-lower-bound} on that subsequence.
This proves \eqref{mf:uniform-positive} and the version with uniform
\(\varepsilon,s\).
\end{proof}

\section{The exact spherical saddles and a stronger Gini bound}\label{mf:sphere}

\subsection{The anisotropic spherical potential}

Let \(Y\) be uniform on the real sphere \(S^2\) and
\[
 v=T_\tau Y,\qquad T_\tau=\operatorname{diag}(1,1,\tau),
 \qquad 0<\tau<1.
\]
The marginal law of \(v\) has covariance
\(\Sigma=\operatorname{diag}(1/3,1/3,\tau^2/3)\).
The order is specified below and does not affect \(\Phi\).

\begin{proposition}[Exact coherent transition]\label{mf:spherical-saddles}
For this spherical law, \(h=0\) is the unique maximizer of \(\Phi\)
when \(t\le3\). When \(t>3\), its maximizers are exactly
\begin{equation}\label{mf:circular-poles}
 h=e^{i\phi}\sqrt{m/2}(1,i,0)
 \quad\hbox{or}\quad
 h=e^{i\phi}\sqrt{m/2}(1,-i,0),
\end{equation}
where \(m\in(0,1)\) is uniquely determined by
\begin{equation}\label{mf:radial-saddle}
 m=m(k):=\int_0^1\frac{k(1-z^2)}{1+k(1-z^2)}\,dz,\qquad
 k=\frac{tm}{2}.
\end{equation}
Equivalently, \(t=2k/m(k)\), a strictly increasing function of \(k>0\),
and
\begin{equation}\label{mf:mass-formula}
 1-m(k)=
 \frac{\operatorname{artanh}\sqrt{k/(1+k)}}{\sqrt{k(1+k)}}.
\end{equation}
\end{proposition}
\begin{proof}
Multiply \(h\) by a phase so that the real and imaginary parts of
\(T_\tau h\) are orthogonal, and rotate \(Y\).
If their squared lengths are \(\lambda_1,\lambda_2\), the logarithmic
term in \(\Phi\) becomes
\[
 \mathbb E\log\bigl(1+t(\lambda_1Y_1^2+\lambda_2Y_2^2)\bigr).
\]
It is symmetric and strictly concave along each line with fixed positive
\(\lambda_1+\lambda_2\): the second derivative is the negative
expectation of a nonzero square proportional to
\((Y_1^2-Y_2^2)^2\).
It is therefore maximized at \(\lambda_1=\lambda_2\), and is strictly
increasing in their sum. That sum is
\(\|T_\tau h\|^2\le\|h\|^2\), with equality only if \(h_3=0\).
Thus, at fixed \(\|h\|^2=m>0\), the maximizing directions are precisely
\eqref{mf:circular-poles}.

The remaining radial function is strictly concave:
\[
 \varphi_t(m)=\int_0^1\log(1+(tm/2)(1-z^2))\,dz-m.
\]
Its derivative at zero is \(t/3-1\). It tends to \(-\infty\) as
\(m\to\infty\), so the asserted unique radial maximum follows.
The stationary equation gives \eqref{mf:radial-saddle}; its right side
is strictly between zero and one.
The function \(m(k)/k\) is strictly decreasing, proving the asserted
monotonicity of \(t(k)\). Integrating
 \((1+k-kz^2)^{-1}\) gives \eqref{mf:mass-formula}.
\end{proof}

The transition at \(t=3\) is a transition of the coherent maximizers.
It does not by itself imply a negative endpoint derivative for every
\(t>3\); the ordered score must still exceed one.

\subsection{A quantitative strict Gini inequality}

Let \(R\) have density \(1/[2(1+r^2)^{3/2}]\), independently of a
uniform angle \(\Theta\). For an integrable real random variable \(Z\),
write \(\mathcal C(Z)=\tfrac12\mathbb E|Z-Z'|\), with an independent
copy \(Z'\). Put \(X=\sqrt2R\cos\Theta\).

\begin{lemma}[A stronger spherical Gini gain]\label{mf:strong-gini}
One has
\begin{equation}\label{mf:gini-gain}
 \mathcal C(X)\ge\frac1{\sqrt2}
             +\frac{\pi^2-8}{32\sqrt2\,\pi}
             >\frac1{\sqrt2}+\frac{13}{1000}.
\end{equation}
\end{lemma}
\begin{proof}
Let \(U=|R|\) and
\(g(a,b)=\mathbb E\max(aU,bU')\) for \(a,b\ge0\).
Independent symmetric signs show that
\(g(a,b)=\mathbb E|aR-bR'|\).
The distribution function of \(U\) is \(u/\sqrt{1+u^2}\), so
\[
 g(a,b)=\int_0^\infty
 \left(1-\frac{s^2}{\sqrt{(s^2+a^2)(s^2+b^2)}}\right)\,ds.
\]
For positive \(a,b\), differentiating the expectation of the maximum
gives
\[
 g_{aa}(a,b)=\frac1b\int_0^\infty
 \frac{u^2\,du}
 {(1+u^2)^{3/2}(1+(a/b)^2u^2)^{3/2}}.
\]
Differentiation is justified locally by an integrable majorant when
\(a,b\) range in a compact subset of \((0,\infty)\).
Homogeneity gives \(g_{ab}=-(a/b)g_{aa}\) and
\(g_{bb}=(a/b)^2g_{aa}\).

Set \(H(x)=g(1+x,1-x)\) and \(a=1+x,b=1-x\).
Substitution \(u=s\sqrt{b/a}\) yields
\[
 H''(x)=\frac4{(ab)^{3/2}}\int_0^\infty
 \frac{s^2\,ds}
 {\bigl((1+s^2)^2+4(1-ab)s^2/(ab)\bigr)^{3/2}}.
\]
Because \(4s^2\le(1+s^2)^2\) and \(0<ab\le1\), the denominator
inside the parentheses is at most \((1+s^2)^2/(ab)\).
Hence
\[
 H''(x)\ge4\int_0^\infty\frac{s^2}{(1+s^2)^3}\,ds
          =\frac\pi4.
\]
Symmetry gives \(H'(0)=0\).
The radial distribution function
\(F_R(r)=\tfrac12+r/(2\sqrt{1+r^2})\) gives
\[
 \mathcal C(R)=\int_{\mathbb R}F_R(r)(1-F_R(r))\,dr
             =\int_{\mathbb R}\frac{dr}{4(1+r^2)}=\frac\pi4,
\]
so \(H(0)=\pi/2\).
Integrating twice and rescaling proves
\begin{equation}\label{mf:uniform-convexity}
 g(a,b)\ge\frac\pi4(a+b)
               +\frac\pi{16}\frac{(a-b)^2}{a+b}.
\end{equation}
Continuity extends this to the boundary, interpreting the final term as
zero at \(a=b=0\). For example, continuity follows from
\(\mathbb E U=1\) and the Lipschitz bound for the maximum.

By symmetry \(X\) has the law of \(SR\), with
\(S=\sqrt2|\cos\Theta|\) independent of \(R\).
Here
\[
 \mathbb ES=\frac{2\sqrt2}{\pi},\qquad
 \mathbb ES^2=1,\qquad 0\le S\le\sqrt2.
\]
Average \eqref{mf:uniform-convexity} over \(S,S'\). Since
\(S+S'\le2\sqrt2\) and
\(\mathbb E(S-S')^2=2\operatorname{Var}(S)\), this gives
\[
 \mathcal C(X)=\tfrac12\mathbb Eg(S,S')
 \ge\frac\pi4\mathbb ES+\frac{\pi}{32\sqrt2}\operatorname{Var}(S)
 =\frac1{\sqrt2}+\frac{\pi^2-8}{32\sqrt2\,\pi}.
\]
For the displayed rational bound, use
\(\pi>157/50\), \(\pi<22/7\), and \(\sqrt2<99/70\).
The resulting rational lower bound for the gain exceeds \(13/1000\).
\end{proof}

\section{A bounded-condition consequence}\label{mf:concrete}

Use spherical coordinates
\[
 Y=(\sqrt{1-z^2}\cos\theta,\sqrt{1-z^2}\sin\theta,z),
\]
where \(z\) is uniform on \([-1,1]\) and \(\theta\) is independently
uniform on \([0,2\pi]\). Order rows by the real quotient
\begin{equation}\label{mf:sort-coordinate}
 X_\tau(Y)=\frac{\sqrt2\,\tau z\cos\theta}{\sqrt{1-z^2}}.
\end{equation}
Its values at the two sphere poles can be assigned arbitrarily, since
those points have measure zero.

At the positive saddle in \eqref{mf:circular-poles}, write \(a=\sqrt m\)
and use the unitary coordinates with first basis vector
\((1,i,0)/\sqrt2\). Dividing a row form by its pole value gives
\[
 c=(1,e^{-2i\theta},
          \sqrt2\,\tau z e^{-i\theta}/\sqrt{1-z^2}).
\]
In these coordinates \(q_h=(p/a)c\), with
\[
 p(z)=\frac{k(1-z^2)}{1+k(1-z^2)}.
\]
Consequently the real part of one exterior component is
\begin{equation}\label{mf:effective-gini}
 C_{\mathrm{eff}}:=\Re(W_{02})
 =\frac1m\cdot\frac12
      \mathbb E\bigl[p(z)p(z')|X_\tau-X'_\tau|\bigr].
\end{equation}
Although the quotient is unbounded, \(p\,c\) is bounded. Equivalently,
\(q_h\) is a continuous bounded function on the original compact sphere.
No logarithmic clipping or radial truncation is needed.

\begin{lemma}[Loss from diagonal dilution]\label{mf:gini-loss}
With \(C_X=\mathcal C(\sqrt2R\cos\Theta)\) as above,
\begin{equation}\label{mf:gini-loss-bound}
 C_{\mathrm{eff}}\ge
 \tau\left(C_X-\frac{\sqrt2}{\sqrt k}
                   -\frac{\log(1+k)}k\right).
\end{equation}
\end{lemma}
\begin{proof}
Let \(C_p=\tfrac12\mathbb E[pp'|X_\tau-X'_\tau|]\).
The inequality \(1-pp'\le(1-p)+(1-p')\), followed by the triangle
inequality, gives
\[
 \tau C_X-C_p\le
 \mathbb E[(1-p)|X_\tau|]+(1-m)\mathbb E|X_\tau|.
\]
Direct integration gives
\[
 \mathbb E|X_\tau|=\frac{2\sqrt2\,\tau}{\pi},\qquad
 \mathbb E[(1-p)|X_\tau|]
 =\frac{2\sqrt2\,\tau}{\pi}
                  \frac{\arctan\sqrt k}{\sqrt k}.
\]
For the second identity substitute \(u=\sqrt{1-z^2}\) into the
radial integral. Also
\[
 1-m=\int_0^1\frac{dz}{1+k(1-z^2)}
 \le\int_0^1\frac{dz}{1+k(1-z)}
 =\frac{\log(1+k)}k.
\]
Now use \(\arctan\sqrt k\le\pi/2\),
\(2\sqrt2/\pi<1\), and \(C_{\mathrm{eff}}=C_p/m\ge C_p\).
\end{proof}

\begin{theorem}[An eventual family with condition number below \(20002\)]
\label{mf:20002}
There is an explicit algebraic sequence of real positive-definite
correlation matrices \(C_N\), defined for every sufficiently large order,
such that
\begin{equation}\label{mf:strong-derivative}
 \limsup_{N\to\infty}
 \frac{D_1(C_N)}{N(N-1)\per C_N}<-\frac1{800}.
\end{equation}
Its off-diagonal entries have modulus at most \(60000/N\), its least
eigenvalue tends to one, and its condition number tends to \(20001\).
In particular its condition number is eventually less than \(20002\).
No explicit sufficient order is asserted in this theorem.
\end{theorem}
\begin{proof}
Take
\[
 \tau=\frac{999}{1000},\qquad t=60000.
\]
The saddle parameter satisfies \(k>25000\).
Indeed \(t(k)=2k/m(k)\) is strictly increasing, and
\[
 1-m(25000)\le\frac{\log25001}{25000}<\frac16
\]
implies \(t(25000)<60000\).
Both loss functions in \eqref{mf:gini-loss-bound} decrease with \(k\).
At \(k=25000\), the bounds
\(\sqrt2<3/2\), \(\sqrt{25000}>158\), and \(\log25001<11\) give
\[
 \frac{\sqrt2}{\sqrt k}+\frac{\log(1+k)}k
 <\frac3{316}+\frac{11}{25000}<\frac1{100}.
\]
For example, \(\log25001<11\) follows from \(e>8/3\) and
\((8/3)^{11}>25001\).
Lemmas~\ref{mf:strong-gini} and \ref{mf:gini-loss} therefore imply
\[
 C_{\mathrm{eff}}>
 \tau\left(\frac1{\sqrt2}+\frac3{1000}\right)
 >\frac1{\sqrt2}+\frac2{1000}.
\]
Using \(\sqrt2>5/4\) shows
\[
 2C_{\mathrm{eff}}^2>1+\frac1{200}.
\]
The two saddle circles are conjugate, so their full score is constant.
Theorem~\ref{mf:criterion}, even using only the first exterior sector,
proves \eqref{mf:strong-derivative} for any ordered discretization
converging to this sphere law.

Here is one deterministic algebraic discretization in every order.
Let \(m=\lfloor\sqrt N\rfloor\), and form the \(m^2\) sphere rows
\[
 z_\ell=\frac{2\ell-1}{m}-1,\quad
 \theta_j=\frac{2\pi j}{m},\quad
 Y_{j,\ell}=
  (\sqrt{1-z_\ell^2}\cos\theta_j,
   \sqrt{1-z_\ell^2}\sin\theta_j,z_\ell),
\]
for \(1\le\ell\le m\), \(0\le j<m\).
Apply \(T_\tau\), add \(N-m^2\) copies of \((1,0,0)\), and sort by
\eqref{mf:sort-coordinate}, breaking ties lexicographically by the
original labels. The grid avoids the two poles and all entries are
algebraic.

These empirical sphere measures converge to the uniform sphere law.
The quotient is continuous away from the two poles and its limiting
distribution has no atoms. Its empirical distribution functions therefore
converge uniformly to its continuous distribution function \(F_X\), and
the largest relative tie multiplicity tends to zero.
The sorted empirical position differs uniformly by \(o(1)\) from
\(F_X(X_\tau(Y))\). Hence the joint ordered law converges to that of
\((F_X(X_\tau(Y)),T_\tau Y)\). Its first marginal is uniform on
\([0,1]\), verifying \eqref{mf:empirical}.

Set \(A_N=I+(60000/N)V_NV_N^T\) and normalize its diagonal:
\[
 C_N=\Delta_N^{-1/2}A_N\Delta_N^{-1/2},\qquad
 \Delta_N=\operatorname{diag}(A_N).
\]
Every permutation monomial, and hence the permanent and its derivative,
is multiplied by the same positive factor. Thus the normalized derivative
ratio is unchanged.
The covariance eigenvalues converge to \(1/3,1/3,\tau^2/3\), so the
condition number of \(A_N\) tends to \(1+60000/3=20001\).
Its least eigenvalue is one for large \(N\).
Since \(\Delta_N=I+O(N^{-1})\) in operator norm and \(\|A_N\|\) is
bounded, \(\|C_N-A_N\|\to0\). This proves the asserted eigenvalue and
condition-number limits. The unit bound on \(\|T_\tau Y\|\), and
diagonal normalization factors at least one, give the off-diagonal bound.
\end{proof}

Before correlation normalization, exactly \(N-3\) eigenvalues equal one
when the rows span \(\mathbb R^3\). Nonconstant diagonal normalization
need not preserve that exact multiplicity; only the stated limiting
spectral conclusions are claimed for \(C_N\).
This theorem concerns large orders only and supplies no explicit
sufficient order; the finite certificate of Section~\ref{sec:certificate}
is a separate construction.

\section{A four-dimensional sphere and condition numbers below 257}
\label{mf:four-sphere}

The second exterior sector can make a decisive contribution. We now use
the complete score \eqref{mf:score} for a four-dimensional real spherical
law. The resulting condition-number bound is much smaller, while the
preceding construction retains its advantage of using three-dimensional
latent rows.

Let \(Y\) be uniform on \(S^3\subset\mathbb R^4\), and set
\[
 v=T_\tau Y,\qquad T_\tau=\operatorname{diag}(1,1,\tau,\tau),
 \qquad 0<\tau<1.
\]
Write
\[
 w_1=Y_1+iY_2,\qquad w_2=Y_3+iY_4,\qquad
 S=|w_1|^2,\qquad b=w_2/w_1.
\]
The variable \(S\) is uniform on \([0,1]\), and the two complex phases
are independent and uniform. These statements follow, for example, by
normalizing two independent standard complex Gaussians.
The ratio \(b\) has density
\begin{equation}\label{mf:complex-cauchy}
 \frac{1}{\pi(1+|b|^2)^2}\,d^2b.
\end{equation}
Indeed \(S=(1+|b|^2)^{-1}\), and changing from the uniform \(S\) to the
radius of \(b\) gives the radial density \(2r/(1+r^2)^2\).
Order the rows by \(x=\Re b\); the exceptional set \(w_1=0\) has
measure zero.

\subsection{Exact saddle geometry and marginal formulas}

The proof of Proposition~\ref{mf:spherical-saddles} applies in real
dimension four: rotational invariance and strict concavity in the two
squared singular values force a maximizing complex direction to be
circular in the first coordinate plane.
For \(t>4\), the maximizers are exactly
\[
 h=e^{i\phi}\sqrt{m/2}(1,i,0,0)
 \quad\hbox{or its conjugate},
\]
where
\begin{equation}\label{mf:four-saddle}
 m=\int_0^1\frac{kS}{1+kS}\,dS
   =1-\frac{\log(1+k)}k,\qquad k=\frac{tm}{2}.
\end{equation}
For completeness, the radial potential is
\(\int_0^1\log(1+(tm/2)S)\,dS-m\); it is strictly concave, has
derivative \(t/4-1\) at zero, and tends to \(-\infty\).
This proves uniqueness of the positive radial saddle.
The function \(t(k)=2k/m(k)\) is strictly increasing by the same
integral argument as before.

Put \(K=1+k\), \(\delta=1-m=\log K/k\), and
\[
 p=\frac{kS}{1+kS}=\frac{k}{K+|b|^2}.
\]
The unweighted marginal density of \(x=\Re b\), and its density
weighted by \(p\), are respectively
\begin{align}
 f(x)&=\frac1{2(1+x^2)^{3/2}},\label{mf:four-f}\\
 f_p(x)&=
 \frac1{2(1+x^2)^{3/2}}
 -\frac1{k\sqrt{1+x^2}}
 +\frac1{k\sqrt{K+x^2}}.\label{mf:four-fp}
\end{align}
To check the second formula, use
\[
 \frac{k}{u^2(u+k)}=\frac1{u^2}-\frac1{ku}
                         +\frac1{k(u+k)}
\]
with \(u=1+|b|^2\), and integrate the imaginary coordinate.
Thus \(f_p\ge0\), \(f-f_p\ge0\), and their masses are \(m\) and
\(\delta\), respectively.

Define the centered prefix mass and first moment
\[
 \mathcal A(x)=\int_{-\infty}^x f_p(u)\,du-\frac m2,\qquad
 \mathcal M(x)=\int_{-\infty}^x u f_p(u)\,du.
\]
Elementary antiderivatives give
\begin{align}
 \mathcal A(x)
 &=\frac{x}{2\sqrt{1+x^2}}
   -\frac{\operatorname{arsinh}x-
          \operatorname{arsinh}(x/\sqrt K)}k,\label{mf:four-prefix}\\
 \mathcal M(x)
 &=-\frac1{2\sqrt{1+x^2}}
   +\frac1{\sqrt{K+x^2}+\sqrt{1+x^2}}.\label{mf:four-moment}
\end{align}
For the second identity, the difference
\((\sqrt{K+x^2}-\sqrt{1+x^2})/k\) is rationalized.
The integration constant is fixed by the zero limit at \(-\infty\).

\subsection{The full score in one-dimensional integrals}

Use the unitary basis
\((1,\pm i,0,0)/\sqrt2\), \((0,0,1,\pm i)/\sqrt2\).
If \(\theta=\arg w_1\) and \(a=e^{-2i\theta}\), the row divided by
its first pole value has coefficients
\begin{equation}\label{mf:four-coefficients}
 c=(1,a,\tau b,\tau a\overline b).
\end{equation}
The phase \(a\) is uniform and independent of \(b\). Conditional on
\(x=\Re b\), the imaginary part of \(b\) has a symmetric law, also
after multiplication by \(p\).
Consequently only coordinates zero and two of the averaged vector
\(q_h=(p/\sqrt m)c\) survive, with densities
\[
 \frac1{\sqrt m}\bigl(f_p(x),0,\tau x f_p(x),0\bigr).
\]
Writing
\begin{equation}\label{mf:four-gamma}
 \Gamma_k=\frac1{2m}\iint|x-y|f_p(x)f_p(y)\,dx\,dy,
\end{equation}
we obtain \(\|W_h\|^2=\tau^2\Gamma_k^2\).
The prefix vector gives
\[
 \|Q_h-2Q_h(x)\|^2
   =\frac4m\bigl(\mathcal A(x)^2+\tau^2\mathcal M(x)^2\bigr).
\]
It follows that the exact limiting score at either saddle circle is
\begin{equation}\label{mf:four-full-score}
 \mathcal R_*=
 2\tau^2\Gamma_k^2+
 \frac8m\int_{\mathbb R}(f-f_p)
         \bigl(\mathcal A^2+\tau^2\mathcal M^2\bigr)\,dx+
 \delta^2.
\end{equation}
All terms are finite. This calculation uses the complete
Theorem~\ref{mf:criterion}; omitting the middle term would not give the
bound below.

\begin{lemma}[Bounds for the two nonconstant sectors]\label{mf:four-bounds}
One has
\begin{equation}\label{mf:four-gamma-lower}
 \Gamma_k\ge
 \frac{\pi/4-2\sqrt K/k}{m}.
\end{equation}
If
\[
 \rho_k=\frac12-\frac\delta2-\frac1{\sqrt K+1}>0,
\]
then
\begin{equation}\label{mf:four-mixed-lower}
 \frac8m\int(f-f_p)(\mathcal A^2+\tau^2\mathcal M^2)\,dx
 \ge\frac{8\tau^2\delta}{m}\rho_k^2.
\end{equation}
\end{lemma}
\begin{proof}
For the probability density \(f\), the absolute-difference potential is
\[
 g(x)=\int|x-y|f(y)\,dy=\sqrt{1+x^2}.
\]
Its derivative is \(2F(x)-1=x/\sqrt{1+x^2}\), and its value at zero
is one, proving the formula. The half mean absolute difference of this
law is \(\pi/4\), by the radial calculation in
Lemma~\ref{mf:strong-gini}.
If \(d=f-f_p\), expansion of \(f_p=f-d\) yields
\[
 \frac12\iint|x-y|f_p(x)f_p(y)\,dx\,dy
 \ge\frac\pi4-\int d(x)g(x)\,dx.
\]
The omitted term is
\(\tfrac12\iint|x-y|d(x)d(y)\,dx\,dy\ge0\).
Using \eqref{mf:four-fp},
\begin{align*}
 \int d(x)g(x)\,dx
 &=\frac1k\int_{\mathbb R}
       \left(1-\frac{\sqrt{1+x^2}}{\sqrt{K+x^2}}\right)\,dx\\
 &\le\frac2k\int_0^\infty
       \left(1-\frac{x}{\sqrt{K+x^2}}\right)\,dx
 =\frac{2\sqrt K}{k}.
\end{align*}
Dividing by \(m\) proves \eqref{mf:four-gamma-lower}.

For the mixed term, compare with the unweighted pair
\[
 \mathcal A_\infty(x)=\frac{x}{2\sqrt{1+x^2}},\qquad
 \mathcal M_\infty(x)=-\frac1{2\sqrt{1+x^2}},
\]
which satisfies \(\mathcal A_\infty^2+\mathcal M_\infty^2=1/4\).
Equations~\eqref{mf:four-prefix}--\eqref{mf:four-moment} give the
uniform estimates
\[
 |\mathcal A-\mathcal A_\infty|\le\delta/2,\qquad
 |\mathcal M-\mathcal M_\infty|\le\frac1{\sqrt K+1}.
\]
The first follows also by integrating the even nonnegative density
\(f-f_p\) of mass \(\delta\) from zero to \(x\).
The triangle inequality in \(\mathbb R^2\) therefore gives
\[
 \mathcal A^2+\tau^2\mathcal M^2
 \ge\tau^2(\mathcal A^2+\mathcal M^2)
 \ge\tau^2\rho_k^2.
\]
Integrating against \(f-f_p\), whose mass is \(\delta\), proves
\eqref{mf:four-mixed-lower}.
\end{proof}

\begin{theorem}[An eventual family with condition number below \(257\)]
\label{mf:257}
There is an explicit algebraic sequence of real positive-definite
correlation matrices \(E_N\), defined for every sufficiently large order,
such that
\begin{equation}\label{mf:four-derivative}
 \lim_{N\to\infty}
 \frac{D_1(E_N)}{N(N-1)\per E_N}<-\frac1{1000}.
\end{equation}
Its off-diagonal entries have modulus at most \(1020/N\), its least
eigenvalue tends to one, and its condition number tends to \(256\).
In particular its condition number is eventually below \(257\).
The unnormalized mean-field perturbation has latent rank four.
This theorem asserts no explicit sufficient order.
\end{theorem}
\begin{proof}
Set \(\tau=999/1000\) and \(t=1020\).
Equation~\eqref{mf:four-saddle} implies \(k<510\).
Since \(m(500)>1-7/500>50/51\), monotonicity of \(t(k)\) gives
\(t(500)<1020\), and therefore \(k>500\).
On this interval,
\[
 6<\log(1+k)<7,\qquad
 \frac6{510}<\delta<\frac7{500},\qquad
 22<\sqrt K<\frac{1131}{50}.
\]
For example, the logarithmic bounds follow from
\(e<11/4\), \((11/4)^6<501\), \(e>8/3\), and
\((8/3)^7>511\). These are strict elementary inequalities.
Using \(\pi>157/50\), Lemma~\ref{mf:four-bounds} gives
\begin{align*}
 \Gamma_k
 &>\frac{157/200-2(1131/50)/500}{1-6/510}
   =\frac{295171}{420000},\\
 \rho_k
 &>\frac12-\frac7{1000}-\frac1{23}
   =\frac{10339}{23000}>0.
\end{align*}
In \eqref{mf:four-full-score}, use \(1/m>1\) and
\(\delta^2\ge\tau^2\delta^2\). The resulting entirely rational bound is
\begin{align}
 \mathcal R_*
 &>\left(\frac{999}{1000}\right)^2
 \left[
 2\left(\frac{295171}{420000}\right)^2
 +8\frac6{510}\left(\frac{10339}{23000}\right)^2
 +\left(\frac6{510}\right)^2
 \right]\notag\\
 &=\frac{1505674035849200604129}
         {1498233800000000000000}
 >1+\frac1{250}.\label{mf:four-rational-margin}
\end{align}
The full score is constant on the two conjugate saddle circles.
Theorem~\ref{mf:criterion} and \eqref{mf:ratio} prove
\eqref{mf:four-derivative}.

An explicit discretization uses \(m=\lfloor N^{1/3}\rfloor\) and the
\(m^3\) rows
\[
 S_\ell=\frac{2\ell-1}{2m},\quad
 \theta_j=\frac{2\pi j}{m},\quad \psi_a=\frac{2\pi a}{m},
\]
\[
 Y_{j,a,\ell}=
 (\sqrt{S_\ell}\cos\theta_j,\sqrt{S_\ell}\sin\theta_j,
  \sqrt{1-S_\ell}\cos\psi_a,\sqrt{1-S_\ell}\sin\psi_a).
\]
Here \(1\le\ell\le m\) and \(0\le j,a<m\).
Apply \(T_\tau\), pad to \(N\) rows with \((1,0,0,0)\), and sort by
\[
 \Re b=\sqrt{\frac{1-S_\ell}{S_\ell}}\cos(\psi_a-\theta_j)
\]
with a fixed lexicographic tie convention. All entries are algebraic
and \(w_1\ne0\) on the grid. The number of added rows is \(O(N^{2/3})\).
The unweighted empirical sphere measures converge to the uniform law.
The limiting ordering coordinate has the continuous density \(f\) in
\eqref{mf:four-f}; the same distribution-function argument used in
Theorem~\ref{mf:20002} verifies ordered empirical convergence.

Set \(A_N=I+(1020/N)V_NV_N^T\) and normalize its diagonal as before.
Its covariance tends to
\(\operatorname{diag}(1/4,1/4,\tau^2/4,\tau^2/4)\), so its condition
number tends to \(1+1020/4=256\).
The row norms are at most one. Diagonal normalization preserves the
derivative ratio, changes the operator by \(o(1)\), and gives the claimed
off-diagonal bound and spectral limits.
\end{proof}

Theorems~\ref{mf:20002} and \ref{mf:257} express a tradeoff: the first
uses three-dimensional latent vectors and has the larger derivative
margin, while the second uses four-dimensional latent vectors and has
the much smaller condition-number bound.
Both are eventual-order conclusions from conventional asymptotic proofs.
Neither claims a finite negative Lean certificate.

